\documentclass[12pt]{article}
\usepackage[utf8]{inputenc}
\usepackage[T1]{fontenc}
\usepackage{hyperref}
\title{Yal\'alag Zapotec Language Resources for K-12 Curriculum Development}
\author{Compiled by Claude Sonnet 5.5 (claude-sonnet-5-5) for Andr\'as Kornai}
\date{}
\begin{document}
\maketitle

\section{Introduction}
Yal\'alag Zapotec (ISO 639-3 zpu; Glottocode yala1267) is an Otomanguean language spoken in Villa Hidalgo Yal\'alag, Oaxaca, Mexico. Sources give about 3,500 speakers (2005), while the survey brief says about 5,500. Searches found some linguistic descriptions and a few texts, but no dictionary and no school materials.

\section{Primary Educational Resources}
No primers or textbooks were found in the searches.

\section{Linguistic and Reference Materials}
\begin{itemize}
\item Wikipedia article on Yal\'alag Zapotec, whose bibliography lists the works below: \url{https://en.wikipedia.org/wiki/Yal%C3%A1lag_Zapotec}
\item Alonso Ortiz, article on positional verbs in Yal\'alag Zapotec, \emph{Ling\"u\'istica Mexicana}: \url{https://linguisticamexicana-amla.colmex.mx/index.php/Linguistica_mexicana/article/download/319/351/1459}
\item L\'opez, \emph{Cuento del conejo y la zorra}, a bilingual text (Colmex repository): \url{https://repositorio.colmex.mx/concern/books/pz50gx934?locale=en}
\item SIL working paper that includes Zapotec data (Zapotec to English): \url{https://icu.repo.nii.ac.jp/record/5377/files/02_Zapotec_ZAPtoENG.pdf}
\item Textile and community information on Yal\'alag: \url{https://mexicantextiles.com/grouppages/groups/zapotecsierra/yalalag/index.htm}
\end{itemize}
The Wikipedia bibliography also lists L\'opez and Newberg, \emph{La conjugaci\'on del verbo zapoteco: zapoteco de Yal\'alag} (SIL, 2005), and Avelino's 2004 UCLA dissertation on Yal\'alag phonetics. I did not see direct links to either.

\section{Religious and Community Resources}
No scripture links were seen in the results.

\section{Digital and Audio Resources}
\begin{itemize}
\item Global Recordings Network, audio: \url{https://globalrecordings.net/en/language/zpu}
\end{itemize}

\section{Curriculum Development Context}
No school use was found. The verb conjugation work, the phonological description and the bilingual text could support orthography and reading-material development. A dictionary would first have to be compiled.

\end{document}